\documentclass[10pt]{article}
\usepackage[T1]{fontenc}
\usepackage[utf8]{inputenc}
\usepackage[french]{babel}
\usepackage{newpxtext,newpxmath}
\usepackage[a4paper,top=1.5cm,bottom=1.35cm,left=1.6cm,right=1.6cm]{geometry}
\usepackage{microtype}
\usepackage{xcolor}
\usepackage{titlesec}
\usepackage[most]{tcolorbox}
\usepackage{enumitem}
\usepackage{array}
\usepackage{booktabs}
\usepackage{fancyhdr}

\definecolor{navy}{HTML}{1F3A5F}
\definecolor{grey}{HTML}{54606E}
\definecolor{green}{HTML}{2C6E49}
\definecolor{amber}{HTML}{9C6A16}
\definecolor{red}{HTML}{A6322E}
\definecolor{plum}{HTML}{7A4EA8}
\definecolor{boxgray}{HTML}{F3F5F8}
\definecolor{boxgreen}{HTML}{F1F6F2}
\definecolor{boxedge}{HTML}{D9E0EA}

\setlength{\parindent}{0pt}
\setlength{\parskip}{2.6pt}
\linespread{1.0}

\pagestyle{fancy}\fancyhf{}
\lhead{\small\color{grey}Scella}
\rhead{\small\color{grey}Classement par potentiel}
\cfoot{\small\color{grey}\thepage}
\renewcommand{\headrulewidth}{0.3pt}

\titleformat{\section}{\color{navy}\large\bfseries}{}{0pt}{}
\titlespacing{\section}{0pt}{9pt}{3pt}
\titleformat{\subsection}{\color{navy}\bfseries}{}{0pt}{}
\titlespacing{\subsection}{0pt}{7pt}{1.5pt}

\newtcolorbox{these}{colback=boxgray,colframe=boxedge,boxrule=0.5pt,arc=2.5pt,
  left=8pt,right=8pt,top=4.5pt,bottom=4.5pt}
\newtcolorbox{cle}{colback=boxgreen,colframe=green!45,boxrule=0.8pt,arc=2.5pt,
  left=8pt,right=8pt,top=5pt,bottom=5pt}

\newcommand{\code}[1]{\texttt{\small #1}}
\newcommand{\best}[1]{{\color{green}\textbf{#1}}}
\newcommand{\sz}[1]{{\small\color{grey}\textsf{#1}}}

\begin{document}

\begin{center}
{\Huge\bfseries\color{navy} Scella}\\[4pt]
{\large\color{grey}\itshape Classement complet des composants par potentiel}\\[4pt]
{\small\color{grey} 5 août 2026 \; · \; panel de 5 juges indépendants, 5 lentilles, une base factuelle commune}
\end{center}
\vspace{2pt}

\begin{these}
\textbf{Le verdict.} \textbf{Le Relais gagne} -- non parce qu'il est le plus rare, mais parce qu'il est
\emph{le seul que personne n'attaque}. Et \textbf{le reçu, qui devrait être premier, se fait battre par
lui-même} : ses deux moitiés ne se touchent jamais, et chacune porte exactement le défaut que l'autre
corrigerait.
\end{these}

\section{1 · Méthode}

Un classement est le genre de jugement où une seule perspective biaise le plus. Cinq juges ont donc reçu
\emph{la même base factuelle} -- l'état vérifié dans le code des treize composants -- et \emph{une
lentille distincte chacun}, sans connaissance des verdicts des autres : la \textbf{maturité} rapportée à
l'effort restant ; l'\textbf{argument d'embauche} chez Kovant ; la \textbf{rareté} pure, maturité
délibérément ignorée ; l'\textbf{adversaire}, un directeur technique hostile en revue diligente ; et la
\textbf{démontrabilité}, c'est-à-dire ce qui tient vingt minutes devant un ingénieur sceptique. Le
composite ci-dessous est la moyenne de leurs rangs. La colonne \emph{écart} mesure l'amplitude entre le
meilleur et le pire rang : un écart faible en haut signale un consensus solide, un écart fort signale un
composant dont la valeur dépend entièrement de la question qu'on pose.

\section{2 · Le classement composite}

\renewcommand{\arraystretch}{1.1}
\begin{center}\small
\begin{tabular}{@{}r>{\raggedright\arraybackslash}p{3.05cm}ccccc c c >{\raggedright\arraybackslash}p{5.5cm}@{}}
\toprule
& \textbf{Composant} & \sz{Mat} & \sz{Kov} & \sz{Rar} & \sz{Adv} & \sz{Dém} & \textbf{Moy} & \sz{Éc.} & \textbf{Ce qui le porte, ce qui le plombe} \\
\midrule
\textbf{1} & \textbf{Le Relais} & 2 & \best{1} & 3 & 3 & \best{1} & \textbf{2,0} & 2 & Original \emph{et} construit \emph{et} servi ; le seul sans caveat \\
\textbf{2} & \textbf{Le reçu, pile A} & \best{1} & 6 & 2 & \best{1} & 6 & \textbf{3,2} & 5 & Scelle le réel ; mais rend \code{verified:true} sur une chaîne tronquée \\
\textbf{3} & \textbf{Le reçu, pile B} & 5 & 3 & \best{1} & 6 & 2 & \textbf{3,4} & 5 & Vérifieur à dépendances vides \emph{testé} ; mais rien à vérifier \\
\textbf{4} & \textbf{Le banc de mesure} & 4 & \best{2} & 4 & 4 & 4 & \textbf{3,6} & 2 & Réfute sa propre thèse avec rigueur ; rien n'est rejouable \\
5 & La Couche 2 servie & 3 & 5 & 10 & \best{2} & 5 & 5,0 & \textbf{8} & Seule mesure en argent réel ; catégorie saturée \\
6 & L'essaim $\Phi$ & 7 & 4 & 5 & 7 & 3 & 5,2 & 4 & Tue sa propre borne en source ; concurrence non pure \\
7 & Le Sandwich & 6 & 7 & 9 & 5 & 7 & 6,8 & 4 & Substrat prêt ; rareté seulement dérivée du reçu \\
8 & L'essaim médiateur & 9 & 9 & 8 & 9 & 10 & 9,0 & 2 & Anti-émergence assumée ; catégorie encombrée, 0 ligne \\
9 & Le Flexible MLP & 13 & 8 & 6 & 10 & 9 & 9,2 & \textbf{7} & Seul actif \emph{déductif} ; jamais flashé, hors dépôt \\
10 & Fil C, la Colonne & 12 & 10 & 7 & 8 & 11 & 9,6 & 5 & Objet visé unique ; rareté empruntée, précondition absente \\
11 & Les nano-robots & 8 & 11 & 13 & 11 & 8 & 10,2 & 5 & Réels ; priors \code{sin(i)} auto-qualifiés « non signifiants » \\
12 & Le méta-essaim & 11 & 12 & 11 & 12 & 12 & 11,6 & \textbf{1} & Volume maximal, valeur réalisée nulle -- \emph{seul consensus total} \\
13 & Le runtime nano & 10 & 13 & 12 & 13 & 13 & 12,2 & 3 & Ré-implémente du résolu, et le ré-implémente faux \\
\bottomrule
\end{tabular}
\end{center}

\section{3 · Les trois découvertes structurelles}

\textbf{Le défaut récurrent n'est pas l'inachèvement, c'est l'auto-attestation.} Le juge adverse a vu ce
qu'une analyse composant par composant ratait nécessairement : \textbf{cinq fois, dans cinq organes
indépendants, la boucle de vérification se referme à l'intérieur du processus qui a produit la chose.}
Le reçu de la pile A recalcule la racine sans jamais la comparer à la racine scellée. Le reçu de la pile
B n'ancre rien, si bien que la racine de référence \emph{voyage dans le dossier} qu'on tend à l'auditeur.
Le méta-essaim \emph{possède} le service qui prouverait à quel moteur on parle, et ce service n'a aucun
appelant. La Couche 2 affiche « préflight frontière OK » alors que la rédaction des données personnelles
est compilée dehors. Le banc gèle ses constantes avant l'expérience mais ne commite aucun script
rejouable, et sa porte de test gold répond « à jour » quel que soit l'état réel. Un audit organe par
organe produit cinq bugs sans lien ; la lentille adverse y voit \textbf{une seule thèse} : tout tient
tant que l'évaluateur est l'auteur, et rien ne survit au passage à un tiers. C'est aussi la bonne
nouvelle, car ce ne sont pas cinq chantiers mais \textbf{un seul principe de réparation} -- déplacer
chaque point de vérification \emph{hors} du producteur. Et corriger l'auto-attestation est exactement ce
que le marché achat-conformité paie ; c'est même à peu près la seule chose qu'il paie.

\textbf{Les quelque 9\,500 tests forment une seule réserve de crédibilité partagée.} C'est la trouvaille
de la lentille de démontrabilité, et elle inverse l'ordre des priorités. Un sceptique qui attrape
\textbf{un seul} faux vert -- \code{verified:true} sur une chaîne tronquée sans clé, « préflight OK »
avec les prompts en clair, un manifeste de vrais participants servi par $\kappa \cdot x$, un balayage de
dérive toujours favorable, une empreinte d'audit forgeable qui repasse verte -- cesse instantanément de
croire aux 320, aux 245, aux 118 et aux 49 autres, \emph{y compris ceux qui sont irréprochables}. Cinq
défauts localisés tiennent ainsi en otage l'ensemble du dépôt. Le corollaire est contre-intuitif et
opérationnellement décisif : le meilleur rendement n'est pas de coder le Sandwich, l'essaim médiateur ou
la Colonne -- qui n'ajouteraient que des lignes à défendre -- mais \textbf{d'éteindre les cinq faux
verts}, dont trois sont petits et déjà outillés. Corollaire de présentation, plus surprenant encore :
les composants à zéro ligne sont aujourd'hui \emph{plus sûrs à montrer} que trois des composants
construits, parce qu'un plan ne peut pas mentir -- à condition d'être présentés comme des plans, jamais
glissés dans le même souffle que le code servi.

\textbf{Le reçu se scinde, et les deux moitiés ne se touchent jamais.} C'est le fait le plus important du
classement. La pile A scelle les missions réellement servies, mais elle \textbf{échoue à la première
sonde du métier} : tronquer, re-vérifier, et obtenir \code{verified:true} sans détenir la moindre clé.
La pile B, elle, tient une revendication extraordinaire et \emph{testée} -- un vérificateur dont
l'arborescence de dépendances est littéralement vide, une recompilation de la source vendorée hors
ligne, puis une re-preuve : « un tiers vérifie sans nous faire confiance » est ici un test qui passe, pas
une diapositive -- mais elle \textbf{n'a rien à vérifier}, aucun producteur ne construisant de graphe
hors des exemples et des tests, et elle \textbf{n'ancre rien}. Chacune porte très exactement le défaut
que l'autre corrigerait.

\section{4 · Là où les juges se battent}

\textbf{La Couche 2, écart 8 -- le plus grand désaccord du panel.} Deuxième pour l'adversaire, qui
observe qu'elle détient \emph{la seule mesure en argent réel} du dépôt ($-$27\,\% de coût à qualité
tenue, 27 runs, vraies clés Claude et Gemini) et que ses deux défauts disqualifiants sont des jours de
travail. Dixième pour la rareté, qui objecte que le routage multi-LLM est la catégorie la plus encombrée
du marché, que $-$27\,\% est le résultat le plus reproductible du secteur, et que sa seule propriété non
banale -- la vérification qui échoue si un appel a été omis -- \textbf{appartient au reçu, pas à elle}.
Les deux ont raison : elle a de la valeur \emph{prouvée} et zéro valeur \emph{rare}.

\textbf{Le banc de mesure, deuxième pour Kovant.} L'argument est le plus fort du panel, et il mérite
d'être cité dans son esprit : Kovant peut embaucher dix ingénieurs capables d'écrire un relais SLM ; il
en trouvera très peu qui \emph{gèlent les constantes avant l'expérience}, appliquent une non-infériorité
appariée et une correction de Holm, puis \textbf{publient que leur propre idée ne paie pas}. Pour une
équipe pre-seed dont toute la thèse repose sur « les petits modèles coordonnés suffisent », l'appareil
capable de dire \emph{quand c'est faux} vaut plus qu'une fonctionnalité de plus -- et c'est le même
appareil qui leur remet le seul résultat positif crédible du dépôt, un modèle de 3,8 milliards de
paramètres à hauteur d'un 12 milliards, obtenu par quelqu'un qui cherchait à le réfuter.

\textbf{Le Flexible MLP, écart 7.} Sixième en rareté pure : c'est le \emph{seul actif déductif} d'un
portefeuille entièrement fait de tests et de mesures, et prouver formellement un invariant sur du code
qui s'adapte sur l'appareil est un front de recherche, pas un exercice. Treizième en délai-valeur, pour
des raisons factuelles et non impressionnistes. La nuance juste, apportée par la lentille de rareté : le
brevet déjà déposé rend \textbf{l'idée} non rare, mais un brevet ne porte pas de preuve -- l'artefact
reste rare même si le mécanisme ne l'est pas.

\textbf{Le méta-essaim, écart 1 -- le seul consensus total.} Unanimement avant-dernier, sur les cinq
lentilles, sans exception. Volume de code maximal, valeur réalisée nulle.

\section{5 · Où je me sépare du composite}

Je ne classe pas la pile A et la pile B comme deux actifs distincts, et voici pourquoi : \textbf{les cinq
juges ont convergé, indépendamment et sans y être invités, sur le même geste}. C'est le signal le plus
fort de tout l'exercice.

\begin{cle}
\textbf{Brancher les reçus de la pile A dans le graphe de la pile B, signer la racine en \emph{ed25519}
-- la bibliothèque est \underline{déjà} dans la crate -- et publier l'ancrage \emph{hors} du dossier.}
Le fait structurel qui rend ce geste court : les empreintes de contenu de la pile A sont \textbf{déjà
sans clé}, donc n'importe qui peut déjà les recalculer ; seuls le HMAC de maillon et le pli de racine
exigent le secret. \sz{Cryptographie : taille S. Garde et rotation des clés : taille M.}
\end{cle}

Ce geste transforme deux raretés aujourd'hui \emph{isolées}, qui s'additionnent, en une seule rareté
\emph{composée}, qui se multiplie : un reçu de mission multi-modèles re-prouvable par un binaire dont
l'auditeur peut lire la totalité de la source. Personne ne détient cela. C'est accessoirement le seul
chemin par lequel le profil de souveraineté visé par la Colonne cesse d'être une idée empruntée.

\textbf{Mon classement réel est donc celui-ci : le reçu \emph{réuni} est premier, loin devant.} Le Relais
est premier \emph{parmi les composants tels qu'ils existent aujourd'hui} -- et c'est cette version-là,
honnêtement, que dit le tableau.

Une correction s'impose au passage. L'analyse de partenariats concluait que « le seul actif réellement
différenciant est le reçu ». Le panel montre que cette phrase est \textbf{vraie de la pile B et fausse
de la pile A} -- or c'est la pile A qui scelle les missions réellement servies. La phrase était trop
grosse d'une moitié.

\section{6 · Ce que ça implique pour Kovant}

L'ordre à tenir est \textbf{l'inverse} de ce que suggérait l'analyse de partenariats. \textbf{Ne pas
ouvrir sur le reçu} : il exige deux caveats immédiats -- la pile A rend \code{verified:true} à un
attaquant \emph{sans clé} qui dispose de l'écriture en base, la pile B n'a aucun producteur.
\textbf{Ouvrir sur le Relais}, qui n'en exige aucun : servi, complet, et sur leur thèse exacte. Puis le
banc, en offrant le résultat positif comme soutien à \emph{leur} thèse et les agrégations réfutées comme
preuve qu'on dira la vérité quand une idée d'essaim ne paie pas. Puis la pile B, en présentant le trou du
producteur comme la moitié que le poste fournit.

Et le retournement le plus fort du dossier : le pire défaut du dépôt -- \emph{le contrôleur du
méta-essaim ne vérifie jamais à quel moteur il parle} -- devient, reformulé en question posée à Kovant,
le meilleur moment de l'entretien.

\begin{cle}
\textbf{« Comment prouvez-vous que l'agent qui a répondu est celui que votre manifeste déclarait ? »}
C'est le problème qu'aucun essaim de petits modèles en conformité ne peut éluder, et il arrive avec la
cicatrice qui prouve qu'on l'a vécu.
\end{cle}

\vspace{3pt}
{\color{grey}\hrule height 0.4pt}
\vspace{2pt}
{\footnotesize\color{grey}Classement produit par un panel de cinq juges indépendants recevant la même
base factuelle -- l'état des composants vérifié dans le code, non dans les documents de conception -- et
une lentille distincte chacun. Les rangs par lentille sont ceux rendus par les juges ; le composite est
la moyenne non pondérée de leurs rangs. Une conclusion antérieure de l'auteur est corrigée au §5.}

\end{document}
